%% kapa: exam=1 round=24 year=2013 period=2 form=A subject=accounting subject_ko=회계학 seq=2 number=2
%% source: Q-Net 감정평가사 기출문제 (cst00309) artlSeq=5073802 file=extracted/1008/5073802/2013년 제24회 감정평가사 1차시험 기출문제/2교시(A형).hwp
%% conversion: kordoc 4.24.51 -> tools/gen1.py (auto); review=proofread-sonnet; flags=none
\begin{question}{2}
포괄손익계산서에 관한 설명으로 옳은 것은?
\begin{choices}
  \item 기업은 예외적인 경우를 제외하고 수익에서 매출원가 및 판매비와관리비(물류원가 등을 포함)를 차감한 영업이익(또는 영업손실)을 포괄손익계산서에 구분하여 표시한다.
  \item 비용을 성격별로 분류하는 기업은 매출원가, 감가상각비, 기타 상각비와 종업원급여비용을 포함하여 비용의 기능에 대한 추가 정보를 공시한다.
  \item 재분류조정은 당기나 과거 기간에 기타포괄손익으로 인식되었다가 당기에 자본잉여금으로 재분류된 금액을 의미한다.
  \item 기타포괄손익의 항목은 세후금액으로 표시할 수 없으며, 관련된 법인세 효과 반영 전 금액으로 표시하고 각 항목들에 관련된 법인세 효과는 단일 금액으로 합산하여 표시한다.
  \item 기업은 재무성과를 설명하는 데 필요하다면 특별항목을 비롯하여 추가항목을 포괄손익계산서에 재량적으로 포함할 수 있으며, 사용된 용어와 항목의 배열도 필요하면 수정할 수 있다.
\end{choices}
\end{question}
\answer{1}
